% Figure area_law_primary_regions.
% Owner: chapter/ch24_peps_area_law_primary_regions.tex (def:area_law_primary_regions).
% Used by: gapped-2d-notes AL App. A 2 (primaries and belts), Fig. 8(a).
% Source: 10-geometry.tex, lines 212--249 and the figure
% geometry:layer-picture, lines 253--291 (D_k, pitch s_k, strips t_k).
% Specialization drawn: s = 4t, r = 2t, shifts a = b = 0; rows run with
% the second coordinate increasing upward, columns with the first.
% Ink-to-index: every dot is one closed fine square of side t. The grey
% tinted comb is the belt: the union of the strips [v_{j,i}, v_{j,i}+t]
% removed from each pitch square (lower and left sides). The red enclosures
% are the open pitch interiors U_j and U_{j'}. Each dashed square is one
% scale-r cell C_{k,z} of D_k that straddles a strip; the blue tinted sets
% are the fragments F_{j,z} = closure(C_{k,z} cap U_j). The label R_j is
% inscribed in the whole-cell fragment and R_{j'} hangs on its fragment's
% hull, both in the layer hue. The two blue fragments inside U_j are
% disconnected but form the single region R_j (all fragments in one pitch
% interior make one primary), so R_j is labelled once: every station beside
% the lower fragment touches the dashed straddling cell or a dot. The
% upper-left fragment of R_j is a whole cell whose dashed outline coincides
% with it and is omitted.
% Colours follow the notes' Section 4.1 table (hue-only redeclarations):
% belt grey (context), layer blue (as D_k in area_law_dyadic_layers).
% Not to scale: the source has t << r << s. No strings, legs or
% contractions; empty boundary signature.
\begin{center}
    \begin{tenkz}[size=s, rows={ket,ket,ket,ket}, cols=8, bonds=none]
        \tndeclare{species}{belt}{hue=source:gray}
        \tndeclare{species}{layer}{hue=source:blue!65!black}
        \tnmark[form=enclosure, species=belt, tint]{(1,1) .. (4,8) - (1,2) .. (3,4) - (1,6) .. (3,8)}{}
        \tnmark[form=enclosure, species=layer, tint]{(1,3) .. (2,4)}{}
        \tnmark[form=enclosure, species=layer, tint]{(3,3) .. (3,4)}{}
        \tnmark[form=enclosure, species=layer, tint, label pos=0]{(1,6) .. (2,6)}{$R_{j'}$}
        \tnmark[form=enclosure, species=complement]{(3,3) .. (4,4)}{}
        \tnmark[form=enclosure, species=complement]{(1,5) .. (2,6)}{}
        \tnmark[form=enclosure, species=selected, label pos=n]{(1,2) .. (3,4)}{$U_j$}
        \tnmark[form=enclosure, species=selected, label pos=n]{(1,6) .. (3,8)}{$U_{j'}$}
        \tnmark[form=label, species=layer, label pos=e]{midway (1,3) and (2,3)}{$R_j$}
    \end{tenkz}
\end{center}
